% =============================================================================
% CHAPTER 19: FINANCIAL DOMAIN MULTI-PERSONA PROMPT GENERATION
% =============================================================================

\chapter{Financial Domain Multi-Persona Prompt Generation}
\label{chap:financial-prompt-gen}

\section{Purpose and Scope}

This chapter defines a generation pipeline for creating synthetic training prompts
in the \textbf{financial domain language}. The pipeline generates diverse, high-quality
natural language prompts that reflect how different \textbf{human populations} and
\textbf{AI agent co-workers} communicate about financial concepts.

\begin{importantbox}
\textbf{Key Distinction:} This chapter generates \textit{prompts and queries} about
financial statements, not the financial data itself. The focus is on the \textbf{language
and communication patterns} of financial professionals and AI agents with varying
levels of domain expertise.
\end{importantbox}

\subsection{Integration Points}

This chapter integrates with:
\begin{itemize}
  \item \textbf{Chapter~\ref{chap:tts-drift}} --- TTS Drift Detection Framework
        (per-persona drift monitoring)
  \item \textbf{Chapter~5 of singapore-ai-safety} --- AI Safety Metrics (fairness,
        explainability, data quality by persona)
  \item \textbf{Chapter~4a of clinerules-agents} --- Agent Spec-Drift Auditing
        (schema expert vs generalist agent drift)
\end{itemize}

\subsection{Core Design Principles}

Following Simula's reasoning-first methodology:

\begin{enumerate}
  \item \textbf{Persona as first-class axis} --- Generation varies by communicator type
  \item \textbf{Expertise-aware templates} --- Schema experts use precise terms; novices are vague
  \item \textbf{Dual-reviewer validation} --- Prompts validated against persona-appropriate standards
  \item \textbf{Controlled diversity} --- All persona $\times$ concept $\times$ intent combinations covered
  \item \textbf{Drift-monitorable} --- Baseline distributions enable per-persona drift detection
\end{enumerate}

\section{Persona Taxonomy}
\label{sec:persona-taxonomy}

\subsection{Persona Classification Axes}

Each communicator is classified along two orthogonal axes:

\begin{table}[htbp]
\centering
\caption{Persona Expertise Axes}
\label{tab:persona-axes}
\begin{tabular}{l l l}
\toprule
\textbf{Axis} & \textbf{Definition} & \textbf{Levels} \\
\midrule
Schema Knowledge & Understanding of database structure, tables, columns & None, Low, Medium, High, Expert \\
Process Knowledge & Understanding of business goals, workflows, outcomes & None, Low, Medium, High, Expert \\
\bottomrule
\end{tabular}
\end{table}

\subsection{Human User Personas}

\begin{table}[htbp]
\centering
\footnotesize
\caption{Human User Persona Definitions}
\label{tab:human-personas}
\begin{tabularx}{\textwidth}{@{}l c c X l@{}}
\toprule
\textbf{Persona} & \textbf{Schema} & \textbf{Process} & \textbf{Communication Style} & \textbf{Code} \\
\midrule
Business Executive & Low & High & High-level, outcome-focused, time-sensitive & HU-EXEC \\
Financial Analyst & Medium & High & Analytical, comparative, ratio-oriented & HU-ANALYST \\
Accountant & High & High & Precise, standards-compliant, audit-ready & HU-ACCT \\
New Hire & Low & Low & Exploratory, uncertain, needs guidance & HU-NEW \\
Operations Manager & Low & Medium & Process-focused, operational metrics & HU-OPS \\
External Auditor & High & Medium & Compliance-focused, evidence-seeking & HU-AUDIT \\
\bottomrule
\end{tabularx}
\end{table}

\subsection{AI Agent Worker Personas}

\begin{table}[htbp]
\centering
\footnotesize
\caption{AI Agent Worker Persona Definitions}
\label{tab:agent-personas}
\begin{tabularx}{\textwidth}{@{}l c c X l@{}}
\toprule
\textbf{Persona} & \textbf{Schema} & \textbf{Process} & \textbf{Communication Style} & \textbf{Code} \\
\midrule
Schema Expert Agent & Expert & Medium & Precise column refs, SQL-like structure & AG-SCHEMA \\
Process Orchestrator & Medium & Expert & Workflow-oriented, task sequencing & AG-ORCH \\
Generalist Helper & Low & High & Natural language, may use synonyms & AG-GEN \\
Data Validation Agent & Expert & Low & Assertion-style, boolean outcomes & AG-VAL \\
Report Generator & Medium & Medium & Template-driven, structured output & AG-REPORT \\
Anomaly Detector & High & Medium & Statistical, threshold-based queries & AG-ANOM \\
\bottomrule
\end{tabularx}
\end{table}

\subsection{Persona Communication Patterns}

\begin{lstlisting}[language=jsonx,caption={Example prompts by persona}]
# HU-EXEC (Business Executive) - Low schema, High process
"What's our revenue growth this quarter compared to last year?"
"Are we on track to hit the margin target?"

# HU-ANALYST (Financial Analyst) - Medium schema, High process  
"Show gross margin by segment with YoY comparison"
"What's the EBITDA trend over the last 4 quarters?"

# HU-ACCT (Accountant) - High schema, High process
"Verify that total assets equals liabilities plus equity for FY2025"
"Reconcile operating income with EBITDA adjustments"

# AG-SCHEMA (Schema Expert Agent) - Expert schema, Medium process
"SELECT SUM(revenue) FROM income_statement WHERE fiscal_year = 2025 
 AND segment_code = 'NAM' GROUP BY quarter"

# AG-GEN (Generalist Helper) - Low schema, High process
"Can you get me the sales numbers for this year?"
"What do our financials look like?"

# HU-NEW (New Hire) - Low schema, Low process
"Where do I find the revenue data?"
"What's a balance sheet?"
\end{lstlisting}

\section{Financial Domain Taxonomy}
\label{sec:financial-taxonomy}

\subsection{Hierarchical Concept Structure}

\begin{lstlisting}[language=jsonx,caption={Financial domain taxonomy (JSON)}]
{
  "version": "1.0",
  "root": "FinancialDomain",
  "nodes": {
    "StatementType": {
      "children": ["IncomeStatement", "BalanceSheet", "CashFlowStatement"]
    },
    "IncomeStatement": {
      "children": ["Revenue", "COGS", "GrossProfit", "OperatingExpenses", 
                   "OperatingIncome", "InterestExpense", "TaxExpense", "NetIncome"],
      "line_items": {
        "Revenue": {"aliases": ["sales", "top line", "turnover"]},
        "COGS": {"aliases": ["cost of sales", "cost of goods"]},
        "GrossProfit": {"aliases": ["gross margin"]},
        "NetIncome": {"aliases": ["bottom line", "profit", "earnings"]}
      }
    },
    "BalanceSheet": {
      "children": ["Assets", "Liabilities", "Equity"],
      "constraints": ["Assets = Liabilities + Equity"]
    },
    "QueryIntent": {
      "children": ["Retrieval", "Comparison", "Validation", "Aggregation", 
                   "Exploration", "Forecasting"]
    },
    "Scenario": {
      "children": ["RoutineReview", "AdhocQuestion", "AuditRequest", 
                   "BoardMeeting", "RegulatoryFiling"]
    }
  }
}
\end{lstlisting}

\subsection{Alias Mapping for Terminology Tolerance}

Schema experts use canonical terms; others may use aliases:

\begin{table}[htbp]
\centering
\footnotesize
\caption{Financial Term Aliases by Persona Expertise}
\label{tab:term-aliases}
\begin{tabularx}{\textwidth}{@{}l l l l@{}}
\toprule
\textbf{Canonical Term} & \textbf{Expert (HU-ACCT)} & \textbf{Analyst (HU-ANALYST)} & \textbf{Novice (HU-NEW)} \\
\midrule
Revenue & Revenue & Revenue, Sales & Sales, Money coming in \\
COGS & Cost of Goods Sold & COGS, Cost of Sales & Costs, What we spend \\
Gross Profit & Gross Profit & Gross Margin & Profit \\
EBITDA & EBITDA & EBITDA & Earnings, Profit \\
Accounts Receivable & A/R, Receivables & Receivables & Money owed to us \\
\bottomrule
\end{tabularx}
\end{table}

\subsection{Data Dictionary and Value Synonym Mapping}

Aliases are not limited to financial concepts. The prompt generator must also
train explicit mappings from surface language to canonical values so that user
phrases are resolved deterministically instead of inferred by the model.

\begin{table}[htbp]
\centering
\footnotesize
\caption{Example Data Dictionary Mappings}
\label{tab:data-dictionary-mappings}
\begin{tabularx}{\textwidth}{@{}l l l X@{}}
\toprule
\textbf{Schema Field} & \textbf{Surface Forms} & \textbf{Canonical Value} & \textbf{Training Requirement} \\
\midrule
\texttt{country\_code} & India, Bharat, Indian entity & \texttt{IN} & Include positive and negative examples for country filters \\
\texttt{country\_code} & Singapore, SG entity, Singapore branch & \texttt{SG} & Include locale-specific prompt variants \\
\texttt{currency\_code} & dollars, US dollars, USD & \texttt{USD} & Disambiguate from SGD, AUD, and local reporting currency \\
\texttt{period} & Q1, first quarter, quarter one & \texttt{Q1} & Include fiscal-calendar context where applicable \\
\texttt{business\_segment} & wealth, WM, wealth management & \texttt{WEALTH} & Require owner-approved segment vocabulary \\
\bottomrule
\end{tabularx}
\end{table}

\noindent
Dictionary entries must be versioned with the training run and linked to schema
paths. Each entry records \texttt{surface\_forms}, \texttt{canonical\_value},
\texttt{locale}, \texttt{owner}, and \texttt{valid\_from}. If a prompt contains a
surface form that is not in the dictionary or schema registry, the generator must
emit a clarification example or reject the training example; it must not invent a
mapping.

\section{Multi-Persona Prompt Generation Pipeline}

\subsection{Stage 1: Mix Generation}

A \textbf{mix} is a concrete combination of:
\begin{itemize}
  \item Persona (from Tables~\ref{tab:human-personas} and~\ref{tab:agent-personas})
  \item Financial concept (from Section~\ref{sec:financial-taxonomy})
  \item Query intent (Retrieval, Comparison, Validation, etc.)
  \item Scenario (Routine, Adhoc, Audit, etc.)
\end{itemize}

\begin{lstlisting}[language=jsonx,caption={Example mix}]
{
  "mix_id": "mix-001",
  "persona": {
    "code": "HU-ANALYST",
    "schema_knowledge": "medium",
    "process_knowledge": "high"
  },
  "financial_concept": "GrossProfit",
  "query_intent": "Comparison",
  "scenario": "RoutineReview",
  "accounting_standard": "US-GAAP"
}
\end{lstlisting}

\subsection{Stage 2: Persona-Aware Meta-Prompt Templates}

Each persona has a distinct template style:

\begin{lstlisting}[language=jsonx,caption={Meta-prompt templates by persona}]
TEMPLATES = {
    "HU-EXEC": """
        Generate a brief, outcome-focused question about {concept}.
        The executive wants to know status or trend, not details.
        Use business language, not technical terms.
        Max 15 words.
    """,
    
    "HU-ANALYST": """
        Generate an analytical query about {concept}.
        Include comparison dimension (YoY, segment, budget vs actual).
        May use financial ratios. Medium technical depth.
        20-40 words.
    """,
    
    "HU-ACCT": """
        Generate a precise accounting query about {concept}.
        Use correct line item names per {standard}.
        May reference specific accounts or reconciliation.
        Include validation or audit perspective.
    """,
    
    "AG-SCHEMA": """
        Generate a structured query about {concept}.
        Use explicit column names and table references.
        May include SQL-like syntax or precise filters.
        Highly technical, unambiguous.
    """,
    
    "AG-GEN": """
        Generate a natural language query about {concept}.
        The agent knows the goal but may not know exact schema terms.
        May use synonyms or vague language.
        Conversational tone.
    """,
    
    "HU-NEW": """
        Generate an exploratory question about {concept}.
        The user is learning and may ask basic questions.
        Uncertain phrasing, may need guidance.
        Simple vocabulary.
    """
}
\end{lstlisting}

\subsection{Stage 3: Complexification by Expertise}

Complexity is \textbf{calibrated to persona expertise}:

\begin{table}[htbp]
\centering
\caption{Complexity Levels by Persona}
\label{tab:complexity-persona}
\begin{tabular}{l c c c c c}
\toprule
\textbf{Persona} & \textbf{c=1} & \textbf{c=2} & \textbf{c=3} & \textbf{c=4} & \textbf{c=5} \\
\midrule
HU-EXEC & \checkmark & \checkmark & -- & -- & -- \\
HU-ANALYST & -- & \checkmark & \checkmark & \checkmark & -- \\
HU-ACCT & -- & -- & \checkmark & \checkmark & \checkmark \\
AG-SCHEMA & -- & -- & \checkmark & \checkmark & \checkmark \\
AG-GEN & \checkmark & \checkmark & \checkmark & -- & -- \\
HU-NEW & \checkmark & \checkmark & -- & -- & -- \\
\bottomrule
\end{tabular}
\end{table}

\textbf{Complexity definitions for financial prompts:}

\begin{table}[htbp]
\centering
\footnotesize
\caption{Complexity Level Definitions}
\label{tab:complexity-levels}
\begin{tabularx}{\textwidth}{@{}c l X@{}}
\toprule
\textbf{Level} & \textbf{Name} & \textbf{Characteristics} \\
\midrule
1 & Basic & Single concept, no aggregation, top-level only \\
2 & Standard & Single concept with filter or time range \\
3 & Detailed & Multi-concept, segment breakdown, ratios \\
4 & Complex & Multi-period comparison, nested aggregation, reconciliation \\
5 & Expert & Multi-statement analysis, adjustments, audit-level detail \\
\bottomrule
\end{tabularx}
\end{table}

\subsection{Stage 4: Dual-Reviewer Validation}

Validation criteria are \textbf{persona-dependent}:

\begin{table}[htbp]
\centering
\footnotesize
\caption{Validation Rules by Persona Type}
\label{tab:validation-persona}
\begin{tabularx}{\textwidth}{@{}l c c c c X@{}}
\toprule
\textbf{Check} & \textbf{HU-EXEC} & \textbf{HU-ANALYST} & \textbf{HU-ACCT} & \textbf{AG-SCHEMA} & \textbf{Rationale} \\
\midrule
Uses canonical terms & Soft & Medium & Strict & Strict & Experts should be precise \\
Schema refs valid & N/A & Soft & Strict & Strict & Schema knowledge varies \\
Ambiguity allowed & Yes & Some & No & No & Novices expected to be vague \\
Complexity appropriate & Low & Med-High & High & High & Match expertise level \\
Grammar correct & Soft & Medium & Strict & N/A & Agents may be terse \\
Intent clear & Medium & Strict & Strict & Strict & All should be answerable \\
\bottomrule
\end{tabularx}
\end{table}

\begin{lstlisting}[language=jsonx,caption={Critic prompt for persona-aware validation}]
def validate_prompt(prompt: str, mix: dict) -> ValidationResult:
    """
    Validate a generated prompt against persona-appropriate criteria.
    """
    persona = mix["persona"]["code"]
    schema_level = mix["persona"]["schema_knowledge"]
    
    # Retrieve persona-specific rules
    rules = VALIDATION_RULES[persona]
    
    violations = []
    
    # Check terminology precision
    if schema_level in ["high", "expert"]:
        if uses_vague_terms(prompt):
            violations.append("terminology_too_vague_for_expert")
    
    # Check complexity appropriateness
    actual_complexity = estimate_complexity(prompt)
    expected_range = COMPLEXITY_RANGES[persona]
    if actual_complexity not in expected_range:
        violations.append(f"complexity_mismatch: {actual_complexity} not in {expected_range}")
    
    # Check ambiguity
    if not rules["ambiguity_allowed"] and is_ambiguous(prompt):
        violations.append("ambiguity_not_allowed_for_persona")
    
    # Check intent clarity
    if not has_clear_intent(prompt):
        violations.append("intent_unclear")
    
    return ValidationResult(
        passed=len(violations) == 0,
        violations=violations,
        persona=persona,
        prompt=prompt,
    )
\end{lstlisting}

\section{Persona-Aware Metrics}
\label{sec:persona-metrics}

\subsection{New Metrics (PER-M01 to PER-M04)}

\begin{table}[htbp]
\centering
\footnotesize
\caption{Persona-Aware Generation Metrics}
\label{tab:persona-metrics}
\begin{tabularx}{\textwidth}{@{}l l X l l@{}}
\toprule
\textbf{ID} & \textbf{Name} & \textbf{Formula} & \textbf{Threshold} & \textbf{Direction} \\
\midrule
\metricid{PER-M01} & Persona Coverage Rate & $\frac{|\text{personas with} \geq N \text{ prompts}|}{|\text{all personas}|}$ & $\geq 95\%$ & Higher better \\
\metricid{PER-M02} & Expertise Consistency & $\frac{|\text{prompts matching persona expertise}|}{|\text{total prompts}|}$ & $\geq 90\%$ & Higher better \\
\metricid{PER-M03} & Cross-Persona Fairness & $\max_i |p_i - \bar{p}|$ where $p_i$ = proportion for persona $i$ & $\leq 10\%$ & Lower better \\
\metricid{PER-M04} & Per-Persona Drift Score & $Z_t^{\text{persona}_i}$ (SSG Z-score per persona) & $Z < 2.0$ & Lower better \\
\bottomrule
\end{tabularx}
\end{table}

\subsubsection{PER-M01: Persona Coverage Rate}

\begin{equation}
\boxed{
PCR = \frac{|\{p : N_p \geq N_{\min}\}|}{|P|}
}
\end{equation}

\begin{description}
  \item[Variables] $P$ = set of all defined personas; $N_p$ = number of prompts for persona $p$;
        $N_{\min}$ = minimum prompts per persona (default 100)
  \item[Threshold] $PCR \geq 0.95$ required
\end{description}

\subsubsection{PER-M02: Expertise Consistency Rate}

\begin{equation}
\boxed{
ECR = \frac{N_{\text{prompts matching expected complexity/terminology}}}{N_{\text{total}}}
}
\end{equation}

\begin{description}
  \item[Definition] A prompt ``matches'' if its complexity and terminology precision
        fall within the expected range for its persona
  \item[Threshold] $ECR \geq 0.90$ required
\end{description}

\subsubsection{PER-M03: Cross-Persona Fairness}

\begin{equation}
\boxed{
CPF = \max_{i \in P} |p_i - \bar{p}|
\quad \text{where} \quad
\bar{p} = \frac{1}{|P|}
}
\end{equation}

\begin{description}
  \item[Definition] Maximum deviation from uniform persona distribution
  \item[Threshold] $CPF \leq 0.10$ required (no persona overrepresented by $>10\%$)
\end{description}

\subsubsection{PER-M04: Per-Persona Drift Score}

Extends the SSG framework (Section~\ref{sec:ssg-framework}) to compute a separate
Z-score for each persona:

\begin{equation}
\boxed{
Z_t^{(p)} = \frac{|x_t^{(p)} - \hat{x}_t^{(p)}|}{\sigma_t^{(p)}}
}
\end{equation}

\begin{description}
  \item[Observable $x_t^{(p)}$] Prompt embedding distance from persona $p$'s training centroid
  \item[Threshold] $Z_t^{(p)} < 2.0$ for each persona (95\% confidence)
  \item[Action] If any persona exceeds threshold, emit persona-specific drift alert
\end{description}

\section{Integration with TTS Drift Framework}
\label{sec:tts-integration}

\subsection{Extended User Population Sampling}

The sampling strategy from Chapter~\ref{chap:tts-drift} is extended with persona weights:

\begin{table}[htbp]
\centering
\footnotesize
\caption{Extended Sampling Strategy by Persona}
\label{tab:extended-sampling}
\begin{tabularx}{\textwidth}{@{}l l l l X@{}}
\toprule
\textbf{Persona Code} & \textbf{Type} & \textbf{Sample Rate} & \textbf{Frequency} & \textbf{Focus Metrics} \\
\midrule
HU-EXEC & Human & 100\% & Real-time & PER-M02, TTS-M09 \\
HU-ANALYST & Human & 50\% & Hourly & TTS-M05, TTS-M07, PER-M02 \\
HU-ACCT & Human & 20\% & Hourly & TTS-M03, TTS-M04, PER-M02 \\
HU-NEW & Human & 100\% & Real-time & TTS-M08 (ARR), PER-M02 \\
AG-SCHEMA & Agent & 100\% & Real-time & TTS-M01, TTS-M17, PER-M04 \\
AG-ORCH & Agent & 50\% & Batch & TTS-M18, PER-M04 \\
AG-GEN & Agent & 20\% & Batch & TTS-M07, TTS-M18, PER-M04 \\
\bottomrule
\end{tabularx}
\end{table}

\subsection{Persona-Attributed Drift Alerts}

Drift alerts from Chapter~\ref{chap:tts-drift} are extended with persona attribution:

\begin{lstlisting}[language=jsonx,caption={Extended drift alert with persona}]
@dataclass
class PersonaDriftAlert(DriftAlert):
    """Drift alert with persona attribution."""
    
    # Inherited fields
    alert_id: str
    severity: str
    drift_type: str
    metric_code: str
    current_value: float
    threshold: float
    
    # Persona extension
    persona_code: str
    persona_type: str  # "human" or "agent"
    schema_knowledge: str
    process_knowledge: str
    
    # Cross-persona context
    other_personas_affected: list[str]
    is_persona_specific: bool  # True if only this persona drifted
    
    def to_audit_trail(self) -> dict:
        return {
            **super().to_dict(),
            "persona": {
                "code": self.persona_code,
                "type": self.persona_type,
                "schema_knowledge": self.schema_knowledge,
                "process_knowledge": self.process_knowledge,
            },
            "cross_persona": {
                "other_affected": self.other_personas_affected,
                "persona_specific": self.is_persona_specific,
            },
        }
\end{lstlisting}

\section{Integration with AI Safety Framework}
\label{sec:ai-safety-integration}

\subsection{Mapping to Singapore AI Safety Metrics}

\begin{table}[htbp]
\centering
\footnotesize
\caption{Persona Metrics Mapped to AI Safety Framework}
\label{tab:ai-safety-mapping}
\begin{tabularx}{\textwidth}{@{}l l X@{}}
\toprule
\textbf{AI Safety Metric} & \textbf{Persona Metric} & \textbf{Application} \\
\midrule
S-01 (Accuracy) & PER-M02 & Prompts accurately reflect persona expertise level \\
S-02 (Robustness) & PER-M04 & Stable performance across persona drift scenarios \\
S-04 (Fairness) & PER-M03 & Equal coverage and quality across personas \\
S-06 (Explainability) & Audit trail & Persona classification rationale is logged \\
S-08 (Data Quality) & PER-M01, M02 & Training data covers all personas appropriately \\
\bottomrule
\end{tabularx}
\end{table}

\subsection{Safety Thresholds by Persona Type}

Human and agent personas have different safety tolerances:

\begin{table}[htbp]
\centering
\caption{Safety Thresholds by Persona Type}
\label{tab:safety-thresholds}
\begin{tabular}{l c c l}
\toprule
\textbf{Metric} & \textbf{Human Threshold} & \textbf{Agent Threshold} & \textbf{Rationale} \\
\midrule
Ambiguity tolerance & Higher & Lower & Humans naturally vague \\
Schema precision & Variable & Strict & Agents should be precise \\
Drift sensitivity & Medium & High & Agent drift more concerning \\
Complexity bounds & Wide & Narrow & Human tasks vary more \\
\bottomrule
\end{tabular}
\end{table}

\section{Orchestration Pipeline}
\label{sec:orchestration}

\subsection{Configuration File}

\begin{lstlisting}[language=yamlx,caption={financial\_prompt\_config.yaml}]
# Financial Domain Prompt Generation Configuration
domain: "Financial professional and agent communication"
taxonomy_version: "1.0"
output_dir: "./generated_prompts"

personas:
  humans:
    - code: HU-EXEC
      schema_knowledge: low
      process_knowledge: high
      target_count: 1000
    - code: HU-ANALYST
      schema_knowledge: medium
      process_knowledge: high
      target_count: 2000
    - code: HU-ACCT
      schema_knowledge: high
      process_knowledge: high
      target_count: 1500
    - code: HU-NEW
      schema_knowledge: low
      process_knowledge: low
      target_count: 500
  agents:
    - code: AG-SCHEMA
      schema_knowledge: expert
      process_knowledge: medium
      target_count: 1500
    - code: AG-ORCH
      schema_knowledge: medium
      process_knowledge: expert
      target_count: 1000
    - code: AG-GEN
      schema_knowledge: low
      process_knowledge: high
      target_count: 1000

global_diversity:
  min_coverage: 0.95
  sampling_strategy: "stratified"

local_diversity:
  prompts_per_mix: 5
  diversity_threshold: 0.85

complexification:
  enabled: true
  persona_aware: true

quality_control:
  critic_model: "gpt-4"
  persona_validation: true
  acceptance_rules:
    critical_failures: ["intent_unclear", "persona_mismatch"]
    warn_on_complexity_violation: true

metrics:
  per_m01_threshold: 0.95
  per_m02_threshold: 0.90
  per_m03_threshold: 0.10
  per_m04_z_threshold: 2.0
\end{lstlisting}

\subsection{Pipeline Implementation}

\begin{lstlisting}[language=jsonx,caption={simula\_financial\_prompt\_generator.py}]
from dataclasses import dataclass, field
from typing import List, Dict, Optional
import itertools

@dataclass
class Persona:
    """A communicator persona with expertise levels."""
    code: str
    type: str  # "human" or "agent"
    schema_knowledge: str
    process_knowledge: str
    target_count: int
    complexity_range: tuple[int, int]

@dataclass
class PromptMix:
    """A combination of factors for prompt generation."""
    mix_id: str
    persona: Persona
    financial_concept: str
    query_intent: str
    scenario: str
    accounting_standard: str = "US-GAAP"

@dataclass
class GeneratedPrompt:
    """A generated prompt with metadata."""
    prompt_id: str
    text: str
    mix: PromptMix
    complexity_level: int
    template_used: str
    validation_result: Optional[dict] = None

class FinancialPromptGenerator:
    """Generate diverse prompts for financial domain communication."""
    
    def __init__(self, config_path: str):
        self.config = load_yaml(config_path)
        self.taxonomy = self._load_taxonomy()
        self.personas = self._load_personas()
        self.teacher = LLMClient(model="gpt-4")
        self.critic = LLMClient(model="gpt-4")
        self.embedder = SentenceTransformer("all-MiniLM-L6-v2")
    
    def run(self) -> List[GeneratedPrompt]:
        """Execute the full generation pipeline."""
        
        # Stage 1: Generate all mixes
        mixes = self._generate_mixes()
        print(f"Generated {len(mixes)} mixes")
        
        # Stage 2: Generate prompts per mix
        all_prompts = []
        for mix in mixes:
            prompts = self._generate_prompts_for_mix(mix)
            all_prompts.extend(prompts)
        print(f"Generated {len(all_prompts)} raw prompts")
        
        # Stage 3: Apply persona-aware complexification
        complexified = []
        for prompt in all_prompts:
            levels = self._get_complexity_levels(prompt.mix.persona)
            for level in levels:
                cp = self._complexify(prompt, level)
                complexified.append(cp)
        print(f"Complexified to {len(complexified)} prompts")
        
        # Stage 4: Validate with persona-aware critic
        validated = []
        for prompt in complexified:
            result = self._validate(prompt)
            if result.passed:
                prompt.validation_result = result.to_dict()
                validated.append(prompt)
        print(f"Validated {len(validated)} prompts")
        
        # Compute metrics
        metrics = self._compute_metrics(validated)
        self._save_artifacts(validated, metrics)
        
        return validated
    
    def _generate_mixes(self) -> List[PromptMix]:
        """Generate Cartesian product of personas x concepts x intents."""
        concepts = self.taxonomy["IncomeStatement"]["children"] + \
                   self.taxonomy["BalanceSheet"]["children"]
        intents = self.taxonomy["QueryIntent"]["children"]
        scenarios = self.taxonomy["Scenario"]["children"]
        
        mixes = []
        for persona in self.personas:
            for concept, intent, scenario in itertools.product(concepts, intents, scenarios):
                mix = PromptMix(
                    mix_id=f"mix-{len(mixes):05d}",
                    persona=persona,
                    financial_concept=concept,
                    query_intent=intent,
                    scenario=scenario,
                )
                mixes.append(mix)
        
        return mixes
    
    def _generate_prompts_for_mix(self, mix: PromptMix) -> List[GeneratedPrompt]:
        """Generate diverse prompts for a single mix."""
        template = TEMPLATES[mix.persona.code]
        base_prompt = template.format(
            concept=mix.financial_concept,
            intent=mix.query_intent,
            scenario=mix.scenario,
            standard=mix.accounting_standard,
        )
        
        # Generate multiple diverse versions
        n = self.config["local_diversity"]["prompts_per_mix"]
        prompts = []
        for i in range(n):
            response = self.teacher.complete(
                f"Generate a unique example for: {base_prompt}"
            )
            prompts.append(GeneratedPrompt(
                prompt_id=f"{mix.mix_id}-{i:02d}",
                text=response,
                mix=mix,
                complexity_level=2,  # Default, will be adjusted
                template_used=mix.persona.code,
            ))
        
        # Deduplicate by embedding similarity
        return self._deduplicate(prompts)
    
    def _validate(self, prompt: GeneratedPrompt) -> ValidationResult:
        """Validate prompt against persona-appropriate criteria."""
        return validate_prompt(prompt.text, prompt.mix.__dict__)
    
    def _compute_metrics(self, prompts: List[GeneratedPrompt]) -> dict:
        """Compute PER-M01 to PER-M04."""
        # PER-M01: Persona Coverage Rate
        counts = {}
        for p in prompts:
            counts[p.mix.persona.code] = counts.get(p.mix.persona.code, 0) + 1
        per_m01 = sum(1 for c in counts.values() if c >= 100) / len(self.personas)
        
        # PER-M02: Expertise Consistency
        consistent = sum(1 for p in prompts if self._is_consistent(p))
        per_m02 = consistent / len(prompts) if prompts else 0
        
        # PER-M03: Cross-Persona Fairness
        proportions = [counts.get(p.code, 0) / len(prompts) for p in self.personas]
        mean_prop = 1 / len(self.personas)
        per_m03 = max(abs(p - mean_prop) for p in proportions)
        
        return {
            "PER-M01": per_m01,
            "PER-M02": per_m02,
            "PER-M03": per_m03,
            "PER-M04": {},  # Computed during drift monitoring
        }
\end{lstlisting}

\section{Implementation Checklist}

\begin{itemize}
  \item[\texttt{[x]}] Define persona taxonomy (human + agent)
  \item[\texttt{[x]}] Define financial domain taxonomy (concepts, intents, scenarios)
  \item[\texttt{[x]}] Create persona-aware prompt templates
  \item[\texttt{[x]}] Define PER-M01 to PER-M04 metrics
  \item[\texttt{[x]}] Specify validation rules by persona
  \item[\texttt{[x]}] Define integration with TTS drift framework
  \item[\texttt{[x]}] Define integration with AI Safety framework
  \item[\texttt{[ ]}] Implement \texttt{simula\_financial\_prompt\_generator.py}
  \item[\texttt{[ ]}] Implement \texttt{simula\_financial\_prompt\_critic.py}
  \item[\texttt{[ ]}] Create JSON schemas for artifacts
  \item[\texttt{[ ]}] Create HANA DDL for persona telemetry
  \item[\texttt{[ ]}] Integrate with CI/CD pipeline
\end{itemize}

\section{Cross-References}

\begin{itemize}
  \item \textbf{TTS Drift Metrics (Chapter~\ref{chap:tts-drift})} --- TTS-M05 (SAS),
        TTS-M07 (TDS), TTS-M08 (ARR), TTS-M18 (PSS) apply to generated prompts
  \item \textbf{SSG Framework (Section~\ref{sec:ssg-framework})} --- Per-persona
        Hull-White calibration for PER-M04
  \item \textbf{AI Safety Metrics (Chapter~5 of singapore-ai-safety)} --- S-01, S-02,
        S-04, S-06, S-08 mapped to persona metrics
  \item \textbf{Spec-Drift Auditing (Chapter~4a of clinerules-agents)} --- DRIFT-009,
        DRIFT-010 for agent persona drift
\end{itemize}
